\documentclass{article}
\usepackage{amsmath}
\usepackage{array}
\usepackage{geometry}
\usepackage{verbatim} 
\usepackage{adjustbox}
\usepackage[table,xcdraw]{xcolor} 
\usepackage{enumitem}
\usepackage{longtable}

\geometry{left=0.5in, right=0.5in, top=0.75in, bottom=0.75in}

%change this to make changes to table format
\setlist[itemize,1]{label=$\circ$, leftmargin=*}

\begin{document}
\section{Literature Study Table}
% %For further references see \href{http://www.overleaf.com}{Something  Linky} or go to the next url: \url{http://www.overleaf.com} or open  the next file \href{run:./file.txt}{File.txt}

% %It's also possible to link directly any word or  \hyperlink{thesentence}{any sentence} in your document.

\small
\begin{adjustbox}{width=\textwidth}
\begin{tabular}{|p{0.11\textwidth}|p{0.11\textwidth}|p{0.11\textwidth}|p{0.11\textwidth}|p{0.11\textwidth}|p{0.11\textwidth}|}
    \hline
    \rowcolor[HTML]{D3D3D3}
    \textbf{Title of Paper} & \textbf{Methods Used} & \textbf{Features} & \textbf{Advantages} & \textbf{Disadvantages}  & \textbf{Future Scope} \\
    \hline
    \vspace{0.1in}
    
    \textbf{A Novel Approach for Crowd Analysis and Density Estimation by using ML Techniques} & 
    \begin{itemize}[leftmargin=*]
        \item MobileNet SSD, 
        \item YOLO v4
        \item Mask-RCNN
    \end{itemize} & 
    \begin{itemize}[leftmargin=*]
        \item Lightweight and versatile
        \item Adaptable to real-time processing
    \end{itemize} & 
    \begin{itemize}[leftmargin=*]
        \item MobileNet SSD: Lightweight, versatile, devices with limited computational power.  
        \item YOLO v4: High precision, real-time processing
        \item Mask-RCNN: Excellent instance segmentation, contextual
    \end{itemize} & 
    \begin{itemize}[leftmargin=*]
    \item MobileNet SSD: Reduced accuracy in high-density crowds.  
    \item YOLO v4: High computational demand.  
    \item Mask-RCNN: Slow processing, computationally heavy
    \end{itemize} & 
   \begin{comment} \begin{itemize}[leftmargin=*]
        \item  Accuracy, precision, recall and F1 score chosen for quantifying detection performance. 
      \textit{\textbf{  \item Average Accuracy :}}
            \item MobileNet SSD: 64.75\%
            \item YoloV4: 84\%
            \item Mask-RCNN: 78.5\%
    \end{itemize} & \end{comment} 
    \begin{itemize}[leftmargin=*]
        \item Real-time optimization for dense crowds
    \end{itemize} \\
    \hline
    \vspace{0.1in}
   \textbf{ A Proposal on Stampede Detection in Real Environments} & 
    \begin{itemize}[leftmargin=*]
        \item Optical flow computation
        \item feature extraction, in real time.
        \item Stacking classifiers to detect stampedes
    \end{itemize} & 
    \begin{itemize}[leftmargin=*]
        \item No thresholds needed
        \item Public dataset use
    \end{itemize} & 
    \begin{itemize}[leftmargin=*]
        \item Real-time detection
        \item Public dataset robustness
        \item Doesn't require predefined thresholds.  
        \item Stacking classifier enhances performance.  
        \item Public datasets ensure robustness.
    \end{itemize} & 
    \begin{itemize}[leftmargin=*]
        \item Limited performance in high-density crowds
            \item Limited performance in high-density crowds.  
        \item Relies on optical flow techniques.
        \item Applicable only to a small number of scenarios.
    \end{itemize} & 
    \begin{comment} \begin{itemize}[leftmargin=*]
        \item Receiver operating characteristic (ROC)
        \item Area under the curve (AUC) metrics 
        \item Cost of Computation
    \end{itemize} &  \end{comment}
     \begin{itemize}[leftmargin=*]
        \item Further research for high-density crowds 
    \end{itemize}\\
    \hline
    \end{tabular}
\newpage

\begin{tabular} {|p{0.13\textwidth}|p{0.14\textwidth}|p{0.13\textwidth}|p{0.15\textwidth}|p{0.13\textwidth}|p{0.13\textwidth}|}
    \hline
    \rowcolor[HTML]{D3D3D3}
    \textbf{Title of Paper} & \textbf{Methods Used} & \textbf{Features} & \textbf{Advantages} & \textbf{Disadvantages} & \textbf{Future Scope} 
    \\
    \hline
     \vspace{0.1in}
      \large \textbf{Background Subtraction in Real Applications: Challenges, Current Models and Future Directions} & 
    \begin{itemize}[leftmargin=*]
        \item {Gaussian Mixture Models (GMM)}: Probabilistic approach to model the background and adapt to changes over time.
        \item ViBe Algorithm
    \end{itemize} & 
    \begin{itemize}[leftmargin=*]
        \item Efficiently separates objects in motion from a static background.
        \item Pixel-based technique that maintains a history of background values and updates dynamically.
        \item Combination with CNN, RNN improve background-foreground segmentation.
    \end{itemize} & 
    \begin{itemize}[leftmargin=*]
        \item GMM provides adaptability to illumination changes and dynamic backgrounds.
        \item Computationally efficient and robust against sudden background variations.
        \item Accuracy in complex environments with occlusions and dynamic elements.
    \end{itemize} & 
    \begin{itemize}[leftmargin=*]
        \item  GMM can be sensitive to sudden lighting changes and requires parameter tuning.
        \item ViBe may struggle with highly dynamic scenes and complex movements.
    \end{itemize} & 
    \begin{itemize}[leftmargin=*]
        \item  Optimization for real-time applications 
        \item  Optimization for minimal computational overhead
    \end{itemize} \\

    \hline
    \vspace{0.1in}
  \large \textbf{ CrowdFormer: An Overlap Patching Vision Transformer} & 
    \begin{itemize}[leftmargin=*]
        \item Weakly-supervised crowd counting model with tranformer-encoder and regression
        \item Vision Transformer (ViT) model for crowd analysis. 
        \item Maintains spatial continuity \& reduces information loss.
    \end{itemize} & 
    \begin{itemize}[leftmargin=*]
         \item Preserves long-term dependencies using overlapping patches
         \item Improves performance on crowd density estimation and localization tasks.
         \item Employs Self Attention Mechanism
    \end{itemize} & 
    \begin{itemize}[leftmargin=*]
        \item Reduces information loss from non-overlapping patches.
        \item Preserves local and global contextual information.
        \item Enhances accuracy while maintaining efficiency.
    \end{itemize} & 
    \begin{itemize}[leftmargin=*]
        \item Increased computational complexity compared to traditional ViTs.
        \item Requires careful hyperparameter tuning for optimal performance.
    \end{itemize} & 
     \begin{itemize}[leftmargin=*]
        \item Application in real-time applications and hybrid models.  
    \end{itemize} 
    \\
    \hline
\end{tabular}
\newpage

\begin{tabular}{|p{0.13\textwidth}|p{0.14\textwidth}|p{0.14\textwidth}|p{0.14\textwidth}|p{0.13\textwidth}|p{0.13\textwidth}|}
    \hline
    \rowcolor[HTML]{D3D3D3}
    \textbf{Title of Paper} & \textbf{Methods Used} & \textbf{Features} & \textbf{Advantages} & \textbf{Disadvantages}  & \textbf{Future Scope} \\
      \hline
      \vspace{0.1in}
   \large \textbf{Pyramid Vision Transformer: A Versatile Backbone for Dense Prediction} & 
    \begin{itemize}[leftmargin=*]
        \item Novel Transformer backbone that incorporates a hierarchical structure
        \item Pyramid structure that progressively reduces the spatial dimensions of feature maps while increasing channel depth
    \end{itemize} & 
    \begin{itemize}[leftmargin=*]
        \item Spatial-reduction attention mechanism to mitigate high memory consumption 
        \item Multi-stage processing to refine feature representations at different scales
        \item PVT outperforms CNN-based backbones while maintaining efficiency
    \end{itemize} & 
    \begin{itemize}[leftmargin=*]
        \item Incorporates a hierarchical feature representation similar to CNNs.
        \item Reduces memory and computational overhead via spatial-reduction attention.
        \item Improves performance on dense prediction tasks such as segmentation and object detection.
    \end{itemize} & 
    \begin{itemize}[leftmargin=*]
         \item Requires substantial computational resources compared to traditional CNNs.
    \item Transformer-based models generally need large-scale training datasets for optimal performance.
    \end{itemize} & 
    \begin{itemize}[leftmargin=*]
        \item Further optimizations are needed to enhance real-time processing capabilities \item Improve efficiency in resource-constrained environments
    \end{itemize} \\
    \hline
    \vspace{0.1in}
    \large \textbf{PVT v2: Improved Baselines with Pyramid Vision Transformer }& 
    \begin{itemize}[leftmargin=*]
        \item Improvement of PVT v1 
        \item Employs Self Attention Mechanism
        \item Vision Transformer (ViT) model for crowd analysis. 
        \item Maintains spatial continuity \& reduces information loss.
    \end{itemize} & 
    \begin{itemize}[leftmargin=*]
         \item Linear spatial reduction attention (LSRA) layer
         \item Overlapping patch embedding (OPE) scheme, 
         \item Convolutional feed-forward network (CFFN)
    \end{itemize} & 
    \begin{itemize}[leftmargin=*]
        \item LSRA significantly lowers computational complexity \& makes model more efficient.
        \item Overlapped patch embedding enhances the extraction of local features, improves the overall representation.
        \item CFFN contributes to robust feature learning without fixed positional embeddings.
    \end{itemize} & 
    \begin{itemize}[leftmargin=*]
        \item Increased computational complexity compared to traditional ViTs.
        \item Requires careful hyperparameter tuning for optimal performance.
    \end{itemize} & 
    \begin{itemize}[leftmargin=*]
        \item Integration with additional local operations or adaptations to other domains.
    \end{itemize} \\
     \hline
\end{tabular}
\newpage

\begin{tabular}{|p{0.13\textwidth}|p{0.13\textwidth}|p{0.15\textwidth}|p{0.13\textwidth}|p{0.14\textwidth}|p{0.13\textwidth}|}
    \hline
    \rowcolor[HTML]{D3D3D3}
    \textbf{Title of Paper} & \textbf{Methods Used} & \textbf{Features} & \textbf{Advantages} & \textbf{Disadvantages} &  \textbf{Future Scope} \\
      \hline
       \vspace{0.1in}
   
    \large \textbf{Transformers in Time Series: A Survey} & 
    \begin{itemize}[leftmargin=*]
        \item Sparse Attention Mechanisms 
        \item Alternative positional encoding schemes tailored for continuous data.
        \item Creation of Hybrid Models modeling.
        \item Uses windowing and memory-efficient attention mechanisms to scale to long time series.
    \end{itemize} & 
    \begin{itemize}[leftmargin=*]
        \item Hybrid modelling with convolutional networks (CNN) improves strength in sequential data
        \item Capable of capturing long-range dependencies in time series data.
        \item Handles multi-modal \& irregular time series data more effectively than traditional models.
    \end{itemize} & 
    \begin{itemize}[leftmargin=*]
        \item Reduced quadratic complexity of standard self-attention. 
        \item Advances in sparse attention make Transformers more computationally efficient for time series applications.
    \end{itemize} & 
    \begin{itemize}[leftmargin=*]
        \item High computational complexity, especially for long sequences.
        \item Requires large-scale training data to generalize well.
        \item Model interpretability remains a challenge compared to traditional statistical methods.
        \item Hyperparameter tuning can be complex.
    \end{itemize} & 
    
    \begin{itemize}[leftmargin=*]
        \item Optimization for real-time applications 
        \item  Optimization for minimal computational overhead
        \item Improve interpretability
    \end{itemize} \\
    \hline
     \vspace{0.1in}
   \large \textbf{ JHU- Crowd++ : Large-Scale Crowd Counting Dataset and Benchmark Method} & 
    \begin{itemize}[leftmargin=*]
        \item Large-scale dataset
        \item Benchmark methods
    \end{itemize} & 
    \begin{itemize}[leftmargin=*]
        \item Diverse scenarios
        \item Rich annotations
    \end{itemize} & 
    \begin{itemize}[leftmargin=*]
        \item Diverse scenarios and environmental conditions.  
        \item Rich annotations at both image and head level.
        \item Distractor images.
    \end{itemize} & 
    \begin{itemize}[leftmargin=*]
        \item Annotation complexity
        \item Limited diversity in distractors.
        \item Limited focus on temporal data.
    \end{itemize} & 
    \begin{itemize}[leftmargin=*]
        \item Exploration of temporal data needed
        \item Generalization to unseen scenarios
        \item Enhancing the dataset diversity and leveraging multi-camera systems.
    \end{itemize} \\
    \hline
\end{tabular}
\end{adjustbox}
\end{document}


\begin{comment}
     Machine Learning for Dense Crowd Direction Prediction Using LSTM & 
    \begin{itemize}[leftmargin=*]
        \item LSTM model predicts movement trajectories based on past positions.
        \item Cone of vision concept for influence detection.
        \item "Unit Angle" quantifies these neighbors. 
    \end{itemize} & 
    \begin{itemize}[leftmargin=*]
        \item Trajectory prediction
        \item Cone of vision concept
    \end{itemize} & 
    \begin{itemize}[leftmargin=*]
        \item Suitable for both structured and unstructured crowds.  
        \item Retains long-term dependencies for more accurate predictions.
        \item Improves trajectory prediction for individuals in crowds.
    \end{itemize} & 
    \begin{itemize}[leftmargin=*]
        \item Requires careful tuning for optimal performance.
        \item Expensive requiring high-quality sequential data.  
    \end{itemize} & 
    \begin{itemize}[leftmargin=*]
        \item Real-world testing required due to complexity.
    \end{itemize} \\\end{comment}



        \begin{comment}
    \vspace{0.1in}
     A CNN-RNN Neural Network Join LSTM for Crowd Counting & 
    \begin{itemize}[leftmargin=*]
        \item CNN for feature extraction \&
        \item LSTM for context relevance to estimate crowd density
    \end{itemize} & 
    \begin{itemize}[leftmargin=*]
        \item Estimates crowd density
    \end{itemize} & 
    \begin{itemize}[leftmargin=*]
        \item High accuracy in dense settings
        \item Robust across densities
    \end{itemize} & 
    \begin{itemize}[leftmargin=*]
        \item Poor performance in sparse crowds
    \end{itemize} & 
    \begin{itemize}[leftmargin=*]
        \item Optimization needed for sparse environments
    \end{itemize} \\
    \hline
    \end{comment}


        \begin{comment}
    Fusion of CCTV Video and Spatial Information for Automated Crowd Congestion Monitoring & 
    \begin{itemize}[leftmargin=*]
        \item CCTV video integration with Crowd mobility graphs (CM graph)
        \item Trajectory is predicted with real-time monitoring.
        \item  Kalman filter predictor is used to develop SORT \& DeepSORT for tracking.
    \end{itemize} & 
    \begin{itemize}[leftmargin=*]
        \item Real-time monitoring
        \item Visualization of congestion
    \end{itemize} & 
    \begin{itemize}[leftmargin=*]
        \item Predictive capabilities  
        \item Real-time monitoring.  
        \item  Effective in visualizing congestion across urban spaces.
    \end{itemize} & 
    \begin{itemize}[leftmargin=*]
        \item Depends on video quality \& computational resources
        \item Challenges with dense crowds
    \end{itemize} & 
    \end{comment}
   \begin{comment} \begin{itemize}[leftmargin=*]
        \itemtracking 
        \item  YoloV7-tiny + SORT: MODA (62.8\%), MOTA (70.8\%), Recall(72.2\%), Precision(98.2\%)
        \item YoloV7-tiny + DeepSORT:  MODA (73.7\%), MOTA (81.2\%), Recall(86.2\%), Precision (94.6\%)
    \end{itemize} &  
    \end{comment}
    \begin{comment}
    Improvements needed for high-density scenarios \\
     \end{comment}


         \begin{comment} 
    Abnormal Crowd Behavior Detection Using Motion Information Images and CNN & 
    \begin{itemize}[leftmargin=*]
        \item Optical flow generate Motion Information Images (MIIs)
        \item CNN classified video frames to detect abnormality
    \end{itemize} & 
    \begin{itemize}[leftmargin=*]
        \item MIIs for motion representation
        \item Abnormality detection
    \end{itemize} & 
    \begin{itemize}[leftmargin=*]
        \item  MIIs allow for better representation of crowd motion.  
        \item Achieves state-of-the-art results.
        \item CNN architecture is simple, 3 layers.
    \end{itemize} & 
    \begin{itemize}[leftmargin=*]
        \item Computationally intensive
        \item Limited to binary classification (normal/abnormal behavior).
        \item MII formation takes a considerable amount of time.
    \end{itemize} & 
    \begin{itemize}[leftmargin=*]
      \textbf{  \item UMN Dataset : }
                \item  MII + CNN: 99.08
                \item MII + ResNet-50: 98.89
           
        \textbf{\item PETS2009 :}
                \item  MII + CNN: 98.39
                \item MII + ResNet-50: 97.04
    \end{itemize} & 
    \begin{itemize}[leftmargin=*]
        \item Explore multi-class detection
    \end{itemize} \\
    \hline
\end{comment}
